% ============================================================
%  Chapter 18: Electrical Transport from Categorical Propagation
%  Source: oxford/publications/sources/current-flux-mechanism.tex
% ============================================================

\epigraph{%
  Ohm's law is a partition counting theorem.%
}{}

\section{Current as Categorical State Propagation}
\label{sec:ch18-current}

In the partition framework, electric current is the rate of propagation of
charged partition boundaries through a conductor.
An electron moves from partition cell to adjacent partition cell at the
partition lag time $\taup$.
The macroscopic current density is:
\begin{equation}
  \mathbf{J} = \frac{e}{\taup} \cdot \frac{n_e}{\delta V},
\end{equation}
where $n_e$ is the number of conduction electrons per volume $\delta V$, and
$1/\taup$ is the cell-crossing rate.

\section{Ohm's Law from Partition Lag}
\label{sec:ch18-ohm}

The electric field $\mathbf{E}$ is the gradient of the partition depth field
(Chapter~\ref{ch:charge}).
An applied field biases the partition lag time:
\begin{equation}
  \taup(\mathbf{E}) = \taup(0) \cdot \frac{1}{1 + eE\delta x/\kB T},
\end{equation}
giving a current proportional to $\mathbf{E}$:
\begin{equation}
  \mathbf{J} = \sigma \mathbf{E}, \qquad
  \sigma = \frac{n_e e^2 \taup}{m_e}.
\end{equation}
This is Ohm's law.
The conductivity $\sigma = n_e e^2 \taup / m_e$ is the Drude formula, derived
here from the partition lag time without assuming a scattering picture.

\section{Kirchhoff's Laws}
\label{sec:ch18-kirchhoff}

Kirchhoff's current law (charge conservation at a node) is the partition
analogue of the fact that partition cells cannot be created or destroyed:
$\sum_i I_i = 0$ at every node.
Kirchhoff's voltage law (sum of partition depth gradients around a loop is zero)
follows from the irrotational condition on $\nabla\Depth$ in the absence of
boundary crossings.

\section{Hall Effect and Magnetoresistance}
\label{sec:ch18-hall}

In a magnetic field, the Lorentz force on the moving partition boundary tilts
the lag-time tensor.
The Hall resistivity is:
\begin{equation}
  \rho_{xy} = \frac{B}{n_e e}.
\end{equation}
This follows from the partition Lagrangian's vector potential term
$\mu\dot{\mathbf{x}} \cdot \Apartition$ when $\Apartition$ is the magnetic
vector potential.

%%  SOURCE: integrate full derivations from
%%          oxford/publications/sources/current-flux-mechanism.tex
%%  ADD: four-probe measurement interpretation, drift-diffusion fundamentals,
%%       quantum Hall effect as partition boundary quantisation.

\bigskip
\begin{tcolorbox}[colback=crownblue!5, colframe=crownblue!60!black,
                  title={\textbf{Chapter 18 Summary}}]
\begin{itemize}[noitemsep]
  \item Electric current is the propagation rate of charged partition
    boundaries at the partition lag time $\taup$.
  \item Ohm's law $\mathbf{J} = \sigma\mathbf{E}$ follows from the
    bias on $\taup$ by the applied field.
  \item Kirchhoff's laws are partition conservation statements.
  \item The Hall effect is the magnetic bias of the partition lag-time tensor.
\end{itemize}
\end{tcolorbox}
